
\section{Fuzzing Background}
    Fuzzing is nowaday a widely used approach for software testing and vulnerability assesment. 
    Its main strenght derives from the fact that could rapidly produce various inputs to be fed into a target executable.
    In general a fuzzer can be divided into the \textbf{Input Generator}, that is responsible of provide inputs to the \textbf{Executor} that execute them.
    During the execution the \textbf{Defect Monitor} check if the input executed ends in a crash of the target.\\
    In particular the Input Generator must choose which input provide to the executor between all the inputs in the \textit{Input Space}.
    A subset of the \textit{Input Space} is represented by the \textit{Valid Input Space}.
    A good Input Generator must choose inputs that are also valid inputs, because only by doing this there is a probability of find a vulnerability\cite{Zhu2022}.
    So the need of a precise input selection and mutation arise, in order to discover as many edge cases as possible.\\
    In order to do so, the Input Generator must gain some informations from the Executor. 
    Depending on the amount of informations that the Input Generator has, fuzzing can be classified into 
    \textit{White-box}, \textit{Black-box} or \textit{Grey-box}.
    In Black-box fuzzing the fuzzer does not have any knowledge about the execution.[add citation of black box fuzzer]
    In White-box fuzzing instead the source code must always be available and technique of Symbolic and Concolic execution are used.[add citation on white box fuzzer].
    In Grey-box fuzzing instead the inputs are choosen using, for example, \textbf{Code Coverage} informations (i.e. number of lines of code visited during execution of the target).
    This approach is called \textit{Coverage-Guided Grey-box fuzzing} and is based on prioritizing the inputs that discover new sections of code.
    The best approach to Code Coverage, in terms of overhead and effectiveness, is \textbf{Edge Coverage} which keeps track of which edges in the \gls{CFG} are taken during execution.
    The coverage informations are obtained through the instrumentation of the target application. 
    Depending on the availability of the source code, the instrumentation could be added in different ways:
    \begin{itemize}
        \item The source code is available : In this case the instrumentation is added at \textit{Compile Time}
        \item The source code is not available : In this case there are two different techniques to instrument the application:
        \begin{itemize}
            \item \textbf{\gls{SBI}} : The instrumentation is added before the target is executed but without having the source code, so relying on Binary Instrumentation techniques
            \item \textbf{\gls{DBI}} : The instrumentation is added during the program execution
        \end{itemize}
    \end{itemize}

\section{Problem Statement}
    Coverage-Guided Grey-box fuzzing is one of the most effective techniques in automatic vulnerability disclosure on software, 
    due to its capability of generating wide amount of inputs based on the feedback obtained through Code Coverage informations 
    that guide the fuzzer towards unexplored regions of code and therefore find hidden vulnerabilities.
    However the effectiveness of this approach strongly depends on the availability of the source code.
    Infact both \gls{SBI} and \gls{DBI} have many limitations.
    The main limitation of \gls{DBI} is represented by huge amount of overheads (higher than 600\% overheads according to \cite{zafl}) that makes the exploration of a substantial part of the valid input space almost impossible.[add citation to paper that state the high overheads]
    Instead \gls{SBI} main difficulty is represented by the difficulty in modifying already compiled binaries.
    There are several problems regarding Binary Rewriting such as Indirect Control Transfers (e.g. \texttt{call rax}), 
    shift in binary sections or instructions that cause the need of rewrite the instructions that reference these sections and so on.\\
    For these reasons the preferable technique nowaday is \textit{Compile time Edge Coverage instrumentation} with fuzzers such as 
    AFL++ \cite{aflplusplus}, LibAFL \cite{libafl} o libFuzzer \cite{libfuzzer}.
    However in practice, for industrial secrecy (i.e. proprietary software) or due to precompiled libraries, the applications are 
    often closed source, and so Compile Time Edge Coverage instrumentation become unusable.\\
    In these scenario binary only fuzzing is the only viable approach.
    Since \gls{DBI} has prohibitive overheads by definition, in this thesis I will address the problem of statically instrument 
    a binary with Edge Coverage instructions in a way that the instrumented binary is compatible with AFL++ .\\
    AFL++ has been choosen as fuzzer due to its popularity and because represent the current state of the art implementing 
    all the state of the art techniques disclosed by researchers in last years \cite{aflplusplus}.\\
    The instrumentation must add as little overheads as possible with respect to compile time overheads and must 
    preserve the correctness of the binary.